\begin{afterword}
\summaryhead

The post continues ``Taboo Your Words.'' To describe a thing without its usual word, the author says, you must visualize it as if seeing it for the first time: a bat becomes a tapering wooden rod, and bases become safe positions in an artificial game. The obstacle is the familiar word, which rushes in and replaces the details. The author ranks this skill among the fundamental ones and links it to pragmatism, reductionism and keeping one's purpose.

Categories throw away information. A mental category that only partly matches a goal lets the goal be lost, as when ``getting an education'' takes the place of learning. The rationalist version of Taboo is played with a voluntary handicap: no synonyms and no newly coined handles, and it is judged by whether it serves its purpose. The instruction is to replace the symbol with the substance. ``The Simple Truth'' is offered as an example for ``truth,'' and Bayes's rule for ``evidence.'' The most important word to taboo is left unnamed.

\responsehead

The post improves on the one before it. ``Taboo Your Words'' gave the method and did not show it working on an open question. This post says more exactly what the method asks: the player bans, by choice, the synonyms and new handles that the card does not list. It applies the point to the method's own name. It gives a test for when the exercise has done its job, whether the result ``fulfilled its purpose.'' And it points to a finished application, ``The Simple Truth.''

That application is described more strongly than its text allows. ``The Simple Truth'' is said to describe truth ``without invoking'' words that include ``believe'' and ``real.'' Its published text uses forms of ``believe'' and ``belief'' 54 times and ``reality'' 15 times, including in the narrator's statement of the thesis. The dialogue does describe reality at a lower level, as whatever determines experimental results, so the method can be seen at work there. The list of avoided words is not accurate.

The post's new step is that a category which only partly matches a goal can fail inside one head the way a proxy measure fails in an institution. It is supported by that one example alone. Two other large claims come without argument: that Taboo is as fundamental as asking ``Why?'', and that trying to define terms like ``evidence'' in words goes in circles.

In short: a fuller statement of the method than ``Taboo Your Words,'' with a test for when it has done its job and a real example, but the example is said to avoid words it uses.
\end{afterword}
